% GENERATED FILE. Edit canonical lessons, then run npm run generate:books.
% canonical-source-hash: fnv1a64:03aed97112e0a44b
% canonical-lessons: HI-C277-vahkarega, HI-R277-what-will-happen

\chapter{What Will Happen}
\label{ch:hi-what-will-happen}
\hlchaptermodality{277}

\begin{chapteropening}
\textbf{By the end of this chapter:} \emph{I can say what will happen, for every person (karūṅgā, karoge, karegā, kareṅge), and what will be (hogā, hogī).}

Complete HI-R277-what-will-happen: say what will happen and what will be.
\end{chapteropening}

\section[\texorpdfstring{\hi{वह} \hi{करेगा} (vah karegā)}{वह करेगा (vah karegā)}]{\hi{वह} \hi{करेगा} (vah karegā) — the future, for every person}

\label{lesson:HI-C277-vahkarega}



\begin{quote}
Say \textquotedblleft{}I will do\textquotedblright{}, then \textquotedblleft{}he will come\textquotedblright{}, \textquotedblleft{}they will come\textquotedblright{} and \textquotedblleft{}you will come\textquotedblright{}. (\emph{karūṅgā}, \emph{āegā}, \emph{āeṅge}, \emph{āoge}.) They are four cells of one grid. This lesson fills in the rest.
\end{quote}

\subsection*{You'll want to know: the whole grid}
\begin{quote}
\textbf{\hi{कल} \hi{मैं} \hi{बाज़ार} \hi{जाऊँगा।}}
\end{quote}

\begin{quote}
\emph{kal maiṁ bāzār jāūṅgā.} — \textquotedblleft{}Tomorrow I will go to the market.\textquotedblright{}
\end{quote}

\noindent
\begin{tabularx}{\linewidth}{@{}>{\raggedright\arraybackslash}X>{\raggedright\arraybackslash}X>{\raggedright\arraybackslash}X@{}}
\toprule
\textbf{who} & \textbf{\hi{करना}} & \textbf{\hi{जाना}} \\
\midrule
\textbf{\hi{मैं}} & \textbf{\hi{करूँगा}} / \textbf{\hi{करूँगी}} & \textbf{\hi{जाऊँगा}} / \textbf{\hi{जाऊँगी}} \\
\textbf{\hi{तुम}} & \textbf{\hi{करोगे}} / \textbf{\hi{करोगी}} & \textbf{\hi{जाओगे}} / \textbf{\hi{जाओगी}} \\
\textbf{\hi{आप}} & \textbf{\hi{करेंगे}} / \textbf{\hi{करेंगी}} & \textbf{\hi{जाएँगे}} / \textbf{\hi{जाएँगी}} \\
\textbf{\hi{वह}} & \textbf{\hi{करेगा}} / \textbf{\hi{करेगी}} & \textbf{\hi{जाएगा}} / \textbf{\hi{जाएगी}} \\
\textbf{\hi{हम}} & \textbf{\hi{करेंगे}} & \textbf{\hi{जाएँगे}} \\
\textbf{\hi{वे}} & \textbf{\hi{करेंगे}} / \textbf{\hi{करेंगी}} & \textbf{\hi{जाएँगे}} / \textbf{\hi{जाएँगी}} \\
\bottomrule
\end{tabularx}

After the stem comes a vowel for the person: \textbf{\hi{ऊँ}} for I, \textbf{\hi{ओ}} for \textbf{\hi{तुम}}, \textbf{\hi{ए}} for he or she, \textbf{\hi{एँ}} for several or for \textbf{\hi{आप}}. Then \textbf{-\hi{गा}}, \textbf{-\hi{गी}} or \textbf{-\hi{गे}}, which follows the gender as \textbf{\hi{था}}, \textbf{\hi{थी}}, \textbf{\hi{थे}} do. A man says \textbf{\hi{करूँगा}}, a woman \textbf{\hi{करूँगी}}.

\textbf{\hi{लेना}} and \textbf{\hi{देना}} shorten the stem: \textbf{\hi{मैं} \hi{लूँगा}} (\emph{maiṁ lūṅgā}), \textbf{\hi{वह} \hi{लेगा}} (\emph{vah legā}), \textbf{\hi{हम} \hi{लेंगे}} (\emph{ham leṅge}); \textbf{\hi{मैं} \hi{दूँगा}} (\emph{maiṁ dūṅgā}), \textbf{\hi{वह} \hi{देगा}} (\emph{vah degā}). \textbf{\hi{पीना}} gives \textbf{\hi{वह} \hi{पिएगा}} (\emph{vah piegā}).

\textquotedblleft{}Will not\textquotedblright{} puts \textbf{\hi{नहीं}} in front: \textbf{\hi{मैं} \hi{नहीं} \hi{जाऊँगा}} (\emph{maiṁ nahīṃ jāūṅgā}).

\subsection*{You'll want to know: \hi{होगा}, will be}
\begin{quote}
\textbf{\hi{कल} \hi{बारिश} \hi{होगी।}}
\end{quote}

\begin{quote}
\emph{kal bāriś hogī.} — \textquotedblleft{}It will rain tomorrow.\textquotedblright{}
\end{quote}

\textbf{\hi{है}} and \textbf{\hi{था}} have a future too: \textbf{\hi{होगा}} (\emph{hogā}), \textbf{\hi{होगी}} (\emph{hogī}), \textbf{\hi{होंगे}} (\emph{hoṅge}). It follows its subject as \textbf{\hi{था}}, \textbf{\hi{थी}}, \textbf{\hi{थे}} do. \textbf{\hi{बारिश}} is feminine, so \textbf{\hi{होगी}}.

\begin{itemize}
\raggedright
  \item \textbf{\hi{वह} \hi{घर} \hi{पर} \hi{होगा}} (\emph{vah ghar par hogā}) — he will be at home
  \item \textbf{\hi{कल} \hi{ठंड} \hi{होगी}} (\emph{kal ṭhaṇḍ hogī}) — it will be cold tomorrow
  \item \textbf{\hi{मेरे} \hi{दोस्त} \hi{वहाँ} \hi{होंगे}} (\emph{mere dost vahā̃ hoṅge}) — my friends will be there
\end{itemize}

\subsection*{Your turn}
Say these aloud:

\begin{itemize}
\raggedright
  \item \textquotedblleft{}Tomorrow I will go to the market\textquotedblright{} — \emph{kal maiṁ bāzār jāūṅgā}
  \item \textquotedblleft{}What will you do?\textquotedblright{}, to a friend — \emph{tum kyā karoge?}
\end{itemize}

\begin{tcolorbox}[breakable,colback=teal!4,colframe=teal!35!black,title={Before you move on}]
\textquotedblleft{}He will do\textquotedblright{}? (\textbf{\hi{करेगा}}.) \textquotedblleft{}We will go\textquotedblright{}? (\textbf{\hi{जाएँगे}}.) \textquotedblleft{}I will take\textquotedblright{}? (\textbf{\hi{लूँगा}}, or \textbf{\hi{लूँगी}}.) \textquotedblleft{}It will rain\textquotedblright{}? (\textbf{\hi{बारिश} \hi{होगी}}: \textbf{\hi{होगी}}, because \textbf{\hi{बारिश}} is feminine.)
\end{tcolorbox}
\section[(review)]{Review — what will happen}

\label{lesson:HI-R277-what-will-happen}



\begin{quote}
Say \textquotedblleft{}Yesterday I went to the market\textquotedblright{}. (\emph{kal maiṁ bāzār gayā}.) Now move it: \textquotedblleft{}Tomorrow I will go to the market\textquotedblright{}. (\emph{kal maiṁ bāzār jāūṅgā}.)
\end{quote}

\subsection*{Your turn}
Say these aloud:

\begin{itemize}
\raggedright
  \item Tomorrow I will go to the market, and then I will come home.
  \item Will you come?, to a friend. Yes, I will come.
  \item The train will come at five o'clock.
  \item It will be cold tomorrow, and it will rain.
  \item Arun and Mira will eat at home.
\end{itemize}

\begin{tcolorbox}[breakable,colback=teal!4,colframe=teal!35!black,title={Before you move on}]
\textquotedblleft{}You will go\textquotedblright{}, to a friend? (\textbf{\hi{जाओगे}}, or \textbf{\hi{जाओगी}}.) \textquotedblleft{}He will give\textquotedblright{}? (\textbf{\hi{देगा}}.) \textquotedblleft{}It will be\textquotedblright{}? (\textbf{\hi{होगा}}, or \textbf{\hi{होगी}} for a feminine thing.) \textquotedblleft{}The train will come\textquotedblright{}? (\textbf{\hi{ट्रेन} \hi{आएगी}}: \textbf{\hi{ट्रेन}} is feminine.)
\end{tcolorbox}
